%%%%%%%%%%%%%%%%%%%%%%%%%%%%%%%%%%%%%%%%%%%%%%%%%%%%%%%%%%%%%%%%%%%%%%
\section{Méthodologie et outils}
%%%%%%%%%%%%%%%%%%%%%%%%%%%%%%%%%%%%%%%%%%%%%%%%%%%%%%%%%%%%%%%%%%%%%%

La première phase suit une approche incrémentale : chaque étape produit
un livrable exploitable avant la suivante. Cela garantit à la fois la
qualité des données et la traçabilité de leur préparation.

\begin{table}[H]
\centering
\begin{tabular}{p{2.8cm}p{3.5cm}p{5.8cm}}
\toprule
\textbf{Étape} & \textbf{Objectif} & \textbf{Sortie attendue} \\
\midrule
Collecte & Enregistrer des phrases et mots parlés par des locuteurs natifs. & Fichiers audio + métadonnées de session. \\
Validation & Contrôler la qualité des enregistrements. & Fichiers validés, doublons et erreurs identifiés. \\
Transcription & Rédiger ou reviser la version textuelle. & Transcriptions fiables et exploitable. \\
Segmentation & Lier le signal audio aux unités de texte. & Découpage en phrases et mots avec bornes temporelles. \\
Annotation & Ajouter des informations utiles à l'analyse. & Métadonnées linguistiques, temporelles et de qualité. \\
Export & Préparer le corpus pour la modélisation. & Jeux d'entraînement, validation et test. \\
\bottomrule
\end{tabular}
\caption{Pipeline de préparation du corpus ALORA}
\end{table}

Les outils retenus restent simples et ouverts : Flutter pour la collecte,
FastAPI pour l'API, PostgreSQL pour les métadonnées, FFmpeg pour le
contrôle audio, ainsi que Git/GitHub pour la gestion du code et de la
documentation. Pour la transcription et la segmentation, Praat et ELAN
restent des références utiles pour une annotation fine et un contrôle
manuel de la qualité. En complément, il est également pertinent
d'explorer MaryTTS et ESPnet2, comme le recommande l'encadrant, afin de
mieux situer le corpus dans des pipelines de traitement automatique de
la parole et d'évaluer l'impact de ces outils sur la préparation des
données.

Le choix final dépendra du niveau de précision recherché, de la charge
de travail de transcription et de la possibilité d'intégrer les outils
dans un workflow reproductible.